% !TEX root = ../main.tex

\section{研究区域与数据}

\subsection{研究区概况}

4个站点均位于美国中西部玉米带内，该区域以玉米和大豆种植为主，分别记为A、B、C和D，自西向东排列顺序为A、B、D、C。四站彼此相互独立，站点中心间距平均约452~km，最大约793~km；各站均为76.8~km~$\times$~76.8~km的正方形区域，边界沿等积投影的坐标轴方向延伸，与经纬线方向存在一定偏差。四站中心坐标依次为A（西经95.9085°、北纬42.8746°）、B（西经91.6892°、北纬42.2382°）、C（西经86.7671°、北纬40.7431°）和D（西经89.2372°、北纬40.1071°）。四站2023年地物分类图见图\ref{fig:study_area}，子图采用相同的范围和显示比例，排列仅反映版面顺序，不表示站点间的实际方位。

图\ref{fig:study_area}中的数据来自美国农业部国家农业统计局（USDA NASS）发布的耕地数据层（Cropland Data Layer，CDL）。基于2023年CDL数据统计结果，4个站点均以玉米和大豆为主要作物类型，两者面积之和均超过农作物总面积的90\%。本文将CDL原始类别重分类为玉米、大豆和其它地类3类，其中除玉米和大豆以外的农作物和非农作物地类均归并为其它地类。四站均以玉米和大豆为主要作物，但地物组成与景观格局存在差异。

玉米和大豆的物候信息来自美国农业部国家农业统计局发布的作物进度报告（Crop Progress Report，CPR）。结合该报告，本文以年积日（DOY）116–266作为主要观测窗口，该窗口覆盖玉米和大豆的主要生育阶段。

\par\addvspace{3pt}
\noindent\begin{minipage}{\columnwidth}
\captionsetup{type=figure}
\centering
\begin{subfigure}{3.8cm}\centering
\includegraphics[width=\linewidth]{study_area/fig_study_area_A.png}
\caption{A站点地物分类图}
\end{subfigure}\hspace{3pt}
\begin{subfigure}{3.8cm}\centering
\includegraphics[width=\linewidth]{study_area/fig_study_area_B.png}
\caption{B站点地物分类图}
\end{subfigure}\par\nointerlineskip\vspace{3pt}
\begin{subfigure}{3.8cm}\centering
\includegraphics[width=\linewidth]{study_area/fig_study_area_C.png}
\caption{C站点地物分类图}
\end{subfigure}\hspace{3pt}
\begin{subfigure}{3.8cm}\centering
\includegraphics[width=\linewidth]{study_area/fig_study_area_D.png}
\caption{D站点地物分类图}
\end{subfigure}\par\nointerlineskip\vspace{3pt}
\definecolor{studyCorn}{HTML}{E5BF34}
\definecolor{studySoy}{HTML}{367451}
\definecolor{studyOther}{HTML}{C8C8C8}
{\songti\zihao{-5}\legbox{studyCorn}\,玉米\quad\legbox{studySoy}\,大豆\quad\legbox{studyOther}\,其它地类}
\caption{四个研究站点2023年地物分类图}
\label{fig:study_area}
\end{minipage}\par\addvspace{3pt}

\subsection{数据获取与预处理}

本文采用的多时相遥感数据源自美国国家航空航天局（NASA）地球数据科学系统提供的协同陆地卫星和哨兵2号数据（Harmonized Landsat and Sentinel-2，HLS）第2版。该产品通过对Landsat-8/9 OLI与Sentinel-2 MSI的观测影像进行辐射定标、大气校正、几何精校正以及网格重采样，实现了两类卫星数据的深度融合。本研究选取其空间分辨率为30~m的多时相地表反射率数据，包含蓝光、绿光、红光、近红外、短波红外1和短波红外2共6个核心波段作为深度学习模型的输入特征。

标签数据采用USDA NASS CropScape门户网站提供的2023年CDL产品。CDL数据提供了美国本土逐像素的年度作物类型分类信息，在本研究中用于模型训练的真值标记与精度检验。研究对原始的CDL分类体系进行重分类，统一划分为玉米（Corn）、大豆（Soybean）和其它地类（Others）3个目标类别。重分类后的标签数据在空间投影和分辨率上与对应的多时相HLS影像进行严格对齐。

在预处理阶段，首先对影像进行镶嵌、云掩膜和裁剪处理。随后，对时间序列采用时序线性插值方法补齐无效观测，并统一4个站点的时间跨度与时间间隔。最终，构建得到空间大小为2560$\times$2560像素，时间分辨率为5天、包含31个时间步、空间分辨率为30~m、光谱波段数为6的数据集。样本构建时，将多时相影像及其对应标签裁剪为128~$\times$~128像素图块，相邻图块无重叠；在此基础上，最终每个站点每年获得400个样本用于后续模型训练与验证。

\subsection{研究区空间特征}

本文基于2023年CDL参考标签，从地物组成和景观格局两个方面比较4个站点的空间特征，为后续站内分类与跨站比较提供研究区背景。

\subsubsection{地物组成}

如图\ref{fig:sec13_composition}所示，4个站点均以玉米和大豆为主要地物，两者面积之和占各站总面积的65.57\%～82.16\%。各站玉米面积占比为39.97\%～45.39\%，大豆占比为23.49\%～36.77\%，站点间大豆占比的变化幅度更大。A站两种作物合计占比最高，其它地类仅占17.84\%；B站大豆占比最低，其它地类占比达到34.43\%。因此，4个站点具有相同的主要作物组成，但两种作物与其它地类的面积配置存在差别。

\newpage
\noindent\begin{minipage}{\columnwidth}
\centering
\includegraphics{study_area/fig_sec13_composition.pdf}
\par\nointerlineskip\vspace{3pt}
\classlegend
\captionof{figure}{四个研究站点2023年地物面积占比}
\label{fig:sec13_composition}
\end{minipage}

\FloatBarrier

\subsubsection{景观格局}

景观指数用于定量描述地物组成、斑块大小与形状以及聚集、邻近等空间配置特征，为比较研究区空间格局提供依据\citep{mcgarigal2023fragstats}。地类面积占比已反映景观组成，本文进一步选用斑块面积（AREA）、形状指数（SHAPE）和斑块数刻画两种作物斑块的尺度、轮廓及数量，以边缘密度（ED）、聚集度指数（AI）和同类最近邻距离（ENN）补充描述类别交界密集程度、同类聚集程度和斑块间邻近状况。

参考斑块为2023年30~m标签中8邻域相连的同类像元，不等同于实际农业田块，单像元斑块和研究区外缘斑块均予保留。单像元斑块占各站相应作物斑块总数的44.31\%～54.64\%，却仅占该作物总面积的0.23\%～0.71\%，因此主要作物面积的分布以面积占比观察。图\ref{fig:sec13_area_shape_joint}对四站两种作物采用相同的六档AREA和四档SHAPE，统计各联合组合占本站该作物总面积的比例，保留大小与形状的对应关系；ENN按完整参考斑块等权统计，每个斑块计一次。

四站具有相同的主要作物背景，同类聚集程度也较接近：两种作物的AI为91.22\%～93.92\%（表\ref{tab:sec13_spatial_metrics}），ENN中位数均为67.08~m，四站玉米的AREA面积加权中位数为116.91～130.95~ha（按面积从小到大累计至50\%确定）。上述数值接近并不表示完整空间分布一致。在这一共同背景下，进一步通过面积—形状联合分布（图\ref{fig:sec13_area_shape_joint}）比较主要作物面积所在斑块的差异，斑块数量、邻近距离和边缘密度作为补充。

图\ref{fig:sec13_area_shape_joint}显示，A站两种作物的面积较集中于少数联合组合，形状相对规则：SHAPE小于2的斑块承载了玉米面积的78.27\%、大豆面积的80.82\%，面积占比最大的组合均为100≤AREA<200~ha、1.5≤SHAPE<2。B站两种作物的分布差异突出：玉米面积占比最大的组合为200≤AREA<500~ha、2≤SHAPE<3，大豆则为10≤AREA<50~ha、1≤SHAPE<1.5；大豆的ENN P90为240.00~m，高于其余三站，反映出中位数未体现的邻近距离差异。

C站的联合分布较分散，形状也较复杂：两种作物斑块数和整站ED均为四站最高，小于10~ha的斑块约占数量的九成，但所占作物面积均不足6\%；AREA≥200~ha且SHAPE≥3的斑块仍分别占玉米、大豆总面积的16.78\%和10.09\%，体现出较多小斑块与大型复杂斑块并存的格局。D站两种作物面积占比前三位的联合组合相同，主要面积—形状分布较接近，整站ED为四站最低的51.76~m/ha。

\begin{table*}[!t]
\begin{minipage}{\textwidth}
\centering
\caption{四个研究站点2023年作物斑块数量、聚集与邻近特征及整站边缘密度}
\label{tab:sec13_spatial_metrics}
\begingroup
\songti\rmfamily\zihao{-5}
\setlength{\tabcolsep}{6pt}
\renewcommand{\arraystretch}{1.2}
\begin{tabular*}{\textwidth}{@{\extracolsep{\fill}}cccccc@{}}
\toprule
站点 & 作物 & 斑块数 & AI（\%） & ENN的P90（m） & 整站ED（m/ha） \\
\midrule
\csname @@input\endcsname tables/sec13_landscape_context_rows.tex
\bottomrule
\end{tabular*}
\par\nointerlineskip\vspace{3pt}
\zihao{-5}
\raggedright 注：P90为第90百分位数，按完整参考斑块等权、相邻顺序值线性插值计算。ENN表示研究区内最近同类斑块的边界像元中心距离，不是田块间净空距离。ED覆盖玉米、大豆及其它地类三类景观，每站列一次。
\par
\endgroup
\end{minipage}
\end{table*}

\begin{figure*}[!t]
\begin{minipage}{\textwidth}
\centering
\captionsetup{type=figure}
% Generated presentation; original bins and v3 plot remain unchanged.
% Center the wider canvas without scaling the square plotting frames.
\noindent\makebox[\linewidth][c]{%
\begin{subfigure}{8.9cm}\centering
\includegraphics{study_area/fig_sec13_area_shape_joint_a.pdf}
\caption{玉米面积—形状联合分布}
\end{subfigure}%
\begin{subfigure}{8.9cm}\centering
\includegraphics{study_area/fig_sec13_area_shape_joint_b.pdf}
\caption{大豆面积—形状联合分布}
\end{subfigure}%
}\par\nointerlineskip\vspace{3pt}
\begin{tikzpicture}[x=1cm,y=1cm,every node/.style={inner sep=0pt,font=\songti\rmfamily\zihao{-5}}]
\path[use as bounding box] (0,0) rectangle (13.5,0.72);
\node[anchor=west] at (0,0.36) {站点};
\definecolor{jointsiteA}{HTML}{0072B2}
\fill[jointsiteA] (1.1000,0.36) circle[radius=0.09cm];
\node[anchor=west] at (1.2700,0.36) {A};
\definecolor{jointsiteB}{HTML}{E69F00}
\fill[jointsiteB] (2.0000,0.36) circle[radius=0.09cm];
\node[anchor=west] at (2.1700,0.36) {B};
\definecolor{jointsiteC}{HTML}{009E73}
\fill[jointsiteC] (2.9000,0.36) circle[radius=0.09cm];
\node[anchor=west] at (3.0700,0.36) {C};
\definecolor{jointsiteD}{HTML}{CC79A7}
\fill[jointsiteD] (3.8000,0.36) circle[radius=0.09cm];
\node[anchor=west] at (3.9700,0.36) {D};
\node[anchor=west] at (5.0,0.36) {作物面积占比};
\definecolor{jointbubblegray}{HTML}{777777}
\fill[jointbubblegray] (7.70000,0.36) circle[radius=0.068688659170cm];
\node[anchor=west] at (7.89869,0.36) {1\%};
\fill[jointbubblegray] (9.50000,0.36) circle[radius=0.153592511188cm];
\node[anchor=west] at (9.78359,0.36) {5\%};
\fill[jointbubblegray] (11.50000,0.36) circle[radius=0.307185022376cm];
\node[anchor=west] at (11.93719,0.36) {20\%};
\end{tikzpicture}


\caption{四个研究站点2023年玉米和大豆的面积—形状联合分布}
\label{fig:sec13_area_shape_joint}
\par\nointerlineskip\vspace{3pt}
{\songti\rmfamily\zihao{-5}\raggedright
注：区间左闭右开，末档包含下限。气泡圆面积与该组合占本站该作物总面积的比例成正比，不表示斑块数量；八个站点—作物组各自合计100\%，全图大小尺度一致。每格A左上、B右上、C左下、D右下，错位仅用于区分站点，不表示档内数值差异。气泡为联合档内斑块汇总，不代表单个真实斑块；分档用于展示，不定义形态类型。\par}
\end{minipage}
\end{figure*}
